\section{综合评价与结论}\label{ch:conclusion}

\subsection{模型比较与选择依据}
不同学习算法并不存在脱离数据分布、损失函数及应用目标的统一最优解。表\ref{tab:model-selection}从模型结构、适用条件及主要限制概括正文讨论的代表性方法。

\begin{table}[htbp]\centering
\caption{代表性监督学习方法的综合比较}\label{tab:model-selection}
\begin{tabularx}{\textwidth}{p{25mm}XX}\toprule
方法 & 主要适用情形 & 主要限制\\\midrule
最小二乘 & 线性趋势、连续响应、系数解释 & 异常值及共线性敏感\\
逻辑回归 & 二分类概率建模、对数几率解释 & 线性得分表达受限；完全分离时估计困难\\
Ridge / Lasso & 系数稳定或稀疏建模 & 依赖特征尺度及惩罚强度选择\\
决策树 & 非线性规则、特征交互 & 单树方差较高，需控制复杂度\\
随机森林 & 通过随机化与聚合降低预测波动 & 解释难度与计算成本随集成规模变化\\
AdaBoost & 逐轮组合弱学习器以改善分类 & 对持续错分样本、标签噪声可能敏感\\
SVM & 间隔分类、核化非线性建模 & 特征缩放、核参数与求解规模影响表现\\\bottomrule
\end{tabularx}
\end{table}

\subsection{理论优势与局限性}
本文使用统一的风险最小化视角关联模型、训练准则与数值算法，并通过最小二乘解析解、梯度下降及中心差分实验检验部分求解过程。公式推导明确相应条件，图表用于说明数学结构，而不是替代统计证据。

局限主要在于现有实验集中于回归和基础优化，尚未在同一真实数据集、统一预处理及重复验证协议下比较全部模型。固定随机种子的多项式结果仅为一次构造实验，不能用于推断普遍优劣。后续可使用重复交叉验证、独立测试集及误差不确定性分析扩展验证范围。

\subsection{结论}
从经验风险到总体风险的差异是模型评价的基本出发点；最小二乘和逻辑回归说明损失与统计假设如何共同定义估计问题；正则化、集成学习及最大间隔方法则从不同角度控制模型复杂度或预测波动。实际建模应同时考虑数学适用条件、求解稳定性、数据泄漏防范及独立评价，避免将训练最优直接等同于泛化最优。
